\chapter{Exchange-Rate Regimes and Currency Unions}\label{ch:b3:exchange-rate-regimes}

A country chooses whether its currency floats, is pegged, or is given up. The
choice decides what its central bank can do, how its economy absorbs shocks, and
whether speculators can attack it. This chapter puts the regimes in the record
since 1870, tests the trilemma that constrains them, models the three kinds of
currency crisis, and asks when a group of countries is better off with one
money --- the question the euro has answered in real time.

\begin{openingchart}{Share of the 18 economies of the Macrohistory Database with a
pegged exchange rate, 1870--2020 (percent; the base country of a fixed system,
such as Britain under gold, counts as pegged).}{Jordà, Schularick and Taylor,
Macrohistory Database (\omterm{def:b3:exchange-rate-regimes:regimes}{peg}); our computation}
\begin{tikzpicture}
\begin{axis}[omtime, width=0.72\linewidth, height=0.4\linewidth, xmin=1865, xmax=2025,
  ymin=0, ymax=105, ylabel={Percent pegged}]
  \addplot[omDef, thick, const plot, mark=none] table[x=year, y=share] {figdata/bachelor-3/14-exchange-rate-regimes/out/pegs.csv};
\end{axis}
\end{tikzpicture}
\end{openingchart}

% figdata: pegs_points
In 1870, 24~\% of these economies were pegged; by 1880, 82~\%, as the \omterm{def:b3:exchange-rate-regimes:gold-standard}{gold standard}
spread. It broke in the 1930s (39~\% in 1932), came back as the Bretton Woods dollar
system (100~\% in 1970), and broke again in the 1970s. Since then about half have
pegged, most of them to the Deutsche Mark and then inside the euro.

\begin{recall}
The Year 1 volume stated the trilemma --- a country cannot have a fixed exchange
rate, free capital movements and its own monetary policy all at once --- and the
Mundell--Fleming model. \Cref{ch:b3:current-account} treated \omterm{def:b3:current-account:sudden-stop}{sudden stops}, and
\cref{ch:b3:exchange-rates} interest parity.
\end{recall}

\section{Regimes in the data}\label{sec:b3:exchange-rate-regimes:regimes}

\begin{definition}[Exchange-rate regime; peg; currency board; dollarisation]\label{def:b3:exchange-rate-regimes:regimes}
\index{exchange-rate regime}\index{peg}\index{currency board}\index{dollarisation}
An \emph{exchange-rate regime} is the rule by which a country's exchange rate is
set. A \emph{peg} fixes it against another currency or a basket: a \emph{soft peg}
can be changed by the government, a \emph{hard peg} is built into law. A
\emph{currency board} is a hard peg in which the monetary base must be fully
backed by reserves of the anchor currency. \emph{Dollarisation} abandons the
national currency for a foreign one.
\end{definition}

\begin{definition}[Gold standard]\label{def:b3:exchange-rate-regimes:gold-standard}
\index{gold standard}
Under the \emph{gold standard}, each central bank converts its currency into gold
at a fixed price on demand. Two currencies with fixed gold prices have a fixed
exchange rate between them, within the cost of shipping gold.
\end{definition}

\section{The trilemma, tested}\label{sec:b3:exchange-rate-regimes:trilemma}

If capital is free and the \omterm{def:b3:exchange-rate-regimes:regimes}{peg} credible, \omterm{def:b3:exchange-rates:interest-parity}{uncovered interest parity} with no
expected depreciation forces the home interest rate to follow the base country's
one for one. A floating country can set its own. The test regresses the yearly
change in the home short rate on the change in the base country's rate.

\begin{omchart}{The trilemma test, 1999--2020: yearly change in the short-term interest
rate against the change in the base country's rate (points), pegs and floats, with
their fitted lines.}{Jordà, Schularick and Taylor, Macrohistory Database (stir,
\omterm{def:b3:exchange-rate-regimes:regimes}{peg}\_type, \omterm{def:b3:exchange-rate-regimes:regimes}{peg}\_base); our computation}
\begin{tikzpicture}
\begin{axis}[omscatter, width=0.6\linewidth, height=0.44\linewidth, xmin=-4.8, xmax=2.8,
  ymin=-6, ymax=4, xlabel={Change in the base rate (points)},
  ylabel={Change in the home rate (points)}, legend pos=north west,
  legend style={font=\footnotesize}]
  \addplot[only marks, omDef, mark=*, mark size=1.2pt] table[x=db, y=d] {figdata/bachelor-3/14-exchange-rate-regimes/out/tri_peg.csv};
  \addplot[only marks, omThm, mark=o, mark size=1.4pt] table[x=db, y=d] {figdata/bachelor-3/14-exchange-rate-regimes/out/tri_float.csv};
  \addplot[omDef, thick, no marks, sharp plot] table[x=db, y=peg] {figdata/bachelor-3/14-exchange-rate-regimes/out/tri_fit.csv};
  \addplot[omThm, thick, dashed, no marks, sharp plot] table[x=db, y=float] {figdata/bachelor-3/14-exchange-rate-regimes/out/tri_fit.csv};
  \legend{Pegs, Floats}
\end{axis}
\end{tikzpicture}
\end{omchart}

% figdata: tri_eras
Since 1999 the \omterm{def:b3:exchange-rate-regimes:regimes}{pegs} --- most of them inside the euro, where the rate is literally
common --- have followed the base with a slope of 0.91 (standard error 0.04), the
floats with 0.45 (0.06). The constraint shows in each era, but in the way the
trilemma says it should, not always as \omterm{def:b3:exchange-rate-regimes:regimes}{pegs} against floats: under Bretton Woods,
\omterm{def:b3:exchange-rate-regimes:regimes}{pegs} followed the base least of all (slope 0.28), because capital controls, the
third corner of the triangle, let them keep their own rates. Over more than 130
years the constraints of the trilemma are largely borne out. % ledger: b3rg:ost.trilemma

\section{Currency crises}\label{sec:b3:exchange-rate-regimes:crises}

\begin{definition}[First-, second- and third-generation crisis models]\label{def:b3:exchange-rate-regimes:crisis-models}
\index{first-generation crisis model}\index{second-generation crisis model}\index{third-generation crisis}
A \emph{first-generation crisis model} explains an attack on a \omterm{def:b3:exchange-rate-regimes:regimes}{peg} by fundamentals
inconsistent with it: domestic credit that grows while the \omterm{def:b3:exchange-rate-regimes:regimes}{peg} holds drains
reserves until speculators take the rest. A \emph{second-generation crisis model}
lets the government choose whether to defend, weighing the \omterm{def:b3:exchange-rate-regimes:regimes}{peg} against its costs,
so that expectations of devaluation can fulfil themselves. A \emph{third-generation
crisis} runs through balance sheets: banks and firms with debts in foreign currency
are ruined by the devaluation, which turns a currency crisis into a banking and
output crisis.
\end{definition}

Take the first-generation model in levels. The money supply is domestic credit
plus reserves, $M = D + R$, and credit grows at $\mu$ a year. Money demand is $M = kS -
c\dot S$: it rises with the price level, which under parity in goods is the exchange
rate $S$ (foreign prices fixed at one), and falls with the expected depreciation $\dot S$,
which under uncovered parity is the interest rate's excess over the foreign one.
Under the \omterm{def:b3:exchange-rate-regimes:regimes}{peg} $\bar S$, $\dot S = 0$ and $M = k\bar S$ is constant, so reserves fall by $\mu$ a year. The
\emph{shadow exchange rate} $\tilde S$ is the rate that would hold if the \omterm{def:b3:exchange-rate-regimes:regimes}{peg} were given up
with no reserves left, $D = k\tilde S - c\dot{\tilde S}$.

\begin{proposition}[The timing of the attack]\label{prop:b3:exchange-rate-regimes:attack}
The shadow rate is $\tilde S_t = D_t/k + c\mu/k^2$. Speculators attack when $\tilde S$ reaches $\bar S$, at
$T = R_0/\mu - c/k$: before reserves would run out, and taking the remaining $c\mu/k$ at
once.
\end{proposition}

\begin{proof}
Try $\tilde S = D/k + a$: then $\dot{\tilde S} = \mu/k$ and $k\tilde S - c\dot{\tilde S} = D + ka - c\mu/k$, which equals
$D$ for $a = c\mu/k^2$. Before $\tilde S$ reaches $\bar S$, an attack would leave the rate below the
\omterm{def:b3:exchange-rate-regimes:regimes}{peg}, a capital loss on the foreign currency bought: no one attacks. After, the
\omterm{def:b3:exchange-rate-regimes:regimes}{peg} would offer a sure gain: speculators would have attacked earlier. So the
attack comes when $\tilde S_T = \bar S$, that is $D_T = k\bar S - c\mu/k$, and reserves are then $R_T = k\bar S -
D_T = c\mu/k$. With $D_T = D_0 + \mu T$ and $R_0 = k\bar S - D_0$, $T = (R_0 - c\mu/k)/\mu$. At $T$ money demand falls
from $k\bar S$ to $k\bar S - c\mu/k$, because the expected depreciation, and so the interest rate,
jumps: the attack is the portfolio shift that matches it.
\end{proof}

\begin{omchart}{The first-generation model with the weekend problem's numbers:
reserves (billions of dalens, left) and the exchange rate with its shadow value
(dalens per kest, right).}{Our simulation (\code{attack.py})}
\begin{tikzpicture}
\begin{axis}[omrecord, width=0.44\linewidth, height=0.36\linewidth, xmin=0, xmax=10,
  ymin=0, ymax=32, xlabel={Years}, ylabel={Reserves}]
  \addplot[omDef, thick, mark=none] table[x=t, y=reserves] {figdata/bachelor-3/14-exchange-rate-regimes/out/attack.csv};
\end{axis}
\begin{axis}[omrecord, at={(0.5\linewidth,0)}, width=0.44\linewidth, height=0.36\linewidth,
  xmin=0, xmax=10, ymin=0.9, ymax=1.3, xlabel={Years}, ylabel={Dalens per kest},
  legend pos=north west, legend style={font=\footnotesize}]
  \addplot[omDef, thick, mark=none] table[x=t, y=rate] {figdata/bachelor-3/14-exchange-rate-regimes/out/attack.csv};
  \addplot[omThm, thick, dashed, mark=none] table[x=t, y=shadow] {figdata/bachelor-3/14-exchange-rate-regimes/out/attack.csv};
  \legend{Actual, Shadow}
\end{axis}
\end{tikzpicture}
\end{omchart}

\begin{model}{The first-generation crisis model}
\assumes A \omterm{def:b3:exchange-rate-regimes:regimes}{peg} defended with reserves, and domestic credit growing to finance
deficits; perfect foresight.
\predicts An attack that ends the \omterm{def:b3:exchange-rate-regimes:regimes}{peg} before reserves run out, when the shadow rate
reaches the \omterm{def:b3:exchange-rate-regimes:regimes}{peg}, with no jump in the exchange rate at the attack.
\tested Kaminsky, Lizondo and Reinhart (1998) asked which indicators, crossing a
threshold, signalled a currency crisis within 24 months in past episodes: the ratio
of broad money to reserves, the model's variable, was among the best, together
with exports, output, equity prices and the real exchange rate, which the model
does not use. % ledger: b3rg:klr.signals
A signal warns of a crisis; it does not date the attack, as perfect foresight would.
\breaks Attacks on \omterm{def:b3:exchange-rate-regimes:regimes}{pegs} with ample reserves and sound budgets, as in 1992.
\end{model}

\begin{record}[Sterling leaves the ERM, 16 September 1992]
Sterling had joined the European exchange-rate mechanism with base rates at 15~\%;
they had fallen to 10~\% in eighteen months. In July 1992 German rates rose, and
holding the parity meant British rates following them up. On 16 September, in
exceptionally turbulent conditions and after heavy official purchases of sterling,
the United Kingdom suspended sterling's membership; Italy stopped supporting the
lira, and the peseta was devalued. % ledger: b3rg:boe.erm
No reserve arithmetic forced this. The question was whether the government
would bear the cost of defence.
\end{record}

\begin{proposition}[Multiple equilibria]\label{prop:b3:exchange-rate-regimes:multiple}
Let a central bank hold reserves of 10 and defend until they are gone. Two
speculators can each sell 6 of the currency at a cost of 0.5; a devaluation by 25~\%
gains each seller $0.25 \times 6 = 1.5$. Both selling and neither selling are Nash
equilibria.
\end{proposition}

\begin{proof}
Payoffs (first speculator, second):
\begin{center}\small
\begin{tabular}{lcc}
 & Sell & Hold \\ \hline
Sell & $1, 1$ & $-0.5, 0$ \\
Hold & $0, -0.5$ & $0, 0$ \\ \hline
\end{tabular}
\end{center}
One seller alone (6) cannot exhaust reserves of 10, so the \omterm{def:b3:exchange-rate-regimes:regimes}{peg} holds and he pays
the cost; together (12) they can. If the other sells, selling pays $1 > 0$; if the
other holds, holding pays $0 > -0.5$. Each strategy is a best response to itself.
\end{proof}

With reserves above 12 no attack can succeed, and below 6 one speculator alone
breaks the \omterm{def:b3:exchange-rate-regimes:regimes}{peg}: only the middle range has two equilibria, and in it a \omterm{def:b3:exchange-rate-regimes:regimes}{peg} that
fundamentals could sustain can fall because markets expect it to. What the
government will bear decides the middle range: high unemployment, which makes
defending with high interest rates painful, widens it. % ledger: b3rg:obs.selffulfil

\begin{record}[The Asian crisis, 1997]
The Thai baht averaged 25.3 per dollar in 1996 and 41.4 in 1998; the Korean won
went from 804 to 1{,}403, the Indonesian rupiah from 2{,}342 to 10{,}014. Thai output fell
by 7.6~\% in 1998, Korean by 4.9~\% and Indonesian by 13.1~\% (World Bank, World
Development Indicators). Output collapses of this size are what the third
generation explains: when banks and firms owe dollars, a devaluation multiplies
their debts in home currency, and the currency crisis becomes a banking crisis
and a slump.
\end{record}

\section{Optimum currency areas}\label{sec:b3:exchange-rate-regimes:oca}

\begin{definition}[Optimum currency area]\label{def:b3:exchange-rate-regimes:oca}
\index{optimum currency area}
An \emph{optimum currency area} is a region whose members gain more from sharing a
currency --- lower transaction costs and no exchange risk on their trade --- than
they lose by giving up their own exchange rate and monetary policy. The loss is
small when members trade much with each other, when their shocks are alike, when
labour moves between them, and when a common budget transfers income to members
hit by bad shocks.
\end{definition}

\begin{omfigure}
\begin{tikzpicture}[econ, x=0.5cm, y=0.5cm]
  \draw[econ axes] (0,0) -- (10.5,0) node[right, font=\scriptsize] {trade integration};
  \draw[econ axes] (0,0) -- (0,8.5) node[above, font=\scriptsize] {correlation of shocks};
  \draw[curve, omDef] (1.5,8) -- (9.5,1) node[right, font=\scriptsize] {OCA line};
  \draw[curve, omThm, dashed] (1,1.5) -- (9.5,6.8) node[right, font=\scriptsize] {integration path};
  \node[font=\scriptsize] at (7.8,6.9) {join};
  \node[font=\scriptsize] at (4.2,0.8) {keep own money};
  \fill (3,2.75) circle (2pt); \fill (6.5,4.9) circle (2pt);
  \draw[->, thick] (3.2,2.9) -- (6.3,4.8);
\end{tikzpicture}
\par\smallskip{\small Above the OCA line, the gains of a common currency exceed the
losses. Integration moves countries along the dashed path: trade raises the
correlation of their cycles.}
\end{omfigure}

\begin{model}{Optimum currency areas}
\assumes Shocks that hit members differently, and wages and prices that adjust
slowly, so that a member hit by a bad shock needs a depreciation it no longer has.
\predicts Unions work among members that trade much, share shocks, move labour or
share a budget; elsewhere, members hit by bad shocks suffer long slumps.
\tested Over thirty years of data for twenty industrialised countries, closer trade
links go with more correlated cycles. % ledger: b3rg:fr.endog
That most of the real appreciation of the euro area's periphery was still there
seven years after its inflows stopped (the chart below) is consistent with the
prediction; a test would compare unions with and without the adjustment margins.
\breaks The criteria are endogenous: countries are more likely to meet them after
integrating than before, so a union can create the conditions it needs.
\end{model}

\section{The record of the euro}\label{sec:b3:exchange-rate-regimes:euro}

\begin{omchart}{Consumer prices in six euro-area economies, 1999--2015 (harmonised
index of consumer prices, annual average, 1999 = 100).}{Eurostat (prc\_hicp\_aind); our
rebasing}
\begin{tikzpicture}
\begin{axis}[omtime, width=0.66\linewidth, height=0.42\linewidth, xmin=1998.5, xmax=2015.5,
  ymin=98, ymax=150, ylabel={1999 = 100}, legend pos=north west, legend columns=2,
  legend style={font=\footnotesize}]
  \addplot[black, thick, mark=none] table[x=year, y=DE] {figdata/bachelor-3/14-exchange-rate-regimes/out/euro_prices.csv};
  \addplot[black, thick, dashed, mark=none] table[x=year, y=FR] {figdata/bachelor-3/14-exchange-rate-regimes/out/euro_prices.csv};
  \addplot[omDef, thick, mark=none] table[x=year, y=EL] {figdata/bachelor-3/14-exchange-rate-regimes/out/euro_prices.csv};
  \addplot[omThm, thick, mark=none] table[x=year, y=ES] {figdata/bachelor-3/14-exchange-rate-regimes/out/euro_prices.csv};
  \addplot[omThm, thick, dashed, mark=none] table[x=year, y=IE] {figdata/bachelor-3/14-exchange-rate-regimes/out/euro_prices.csv};
  \addplot[omDef, thick, dotted, mark=none] table[x=year, y=PT] {figdata/bachelor-3/14-exchange-rate-regimes/out/euro_prices.csv};
  \legend{Germany, France, Greece, Spain, Ireland, Portugal}
\end{axis}
\end{tikzpicture}
\end{omchart}

% figdata: euro_prices_points
From 1999 to 2008, consumer prices rose by 17~\% in Germany and by 30 to 36~\% in Greece,
Spain, Ireland and Portugal. With a common currency, that is a real appreciation of
the periphery of 11 to 16~\% against Germany: its goods became dearer while capital flowed in
(\cref{ch:b3:current-account}). When the inflows stopped, the periphery needed
the real depreciation back, without a currency to depreciate.

\begin{definition}[Internal devaluation]\label{def:b3:exchange-rate-regimes:internal-devaluation}
\index{internal devaluation}
An \emph{internal devaluation} is a real depreciation achieved inside a fixed
exchange rate or a currency union by a fall of domestic wages and prices relative
to those of trading partners.
\end{definition}

From 2008 to 2015 Irish prices rose by 0.5~\% (135.7 to 136.4) while German ones rose
by 9.4~\% (117.5 to 128.5): an \omterm{def:b3:exchange-rate-regimes:internal-devaluation}{internal devaluation} of about 8~\% in seven years. Greek,
Spanish and Portuguese prices rose by 7 to 8~\%, barely less than German ones; most of
the periphery's real appreciation was still there in 2015. \omterm{def:b3:exchange-rate-regimes:internal-devaluation}{Internal devaluation}
works through unemployment, which lowers wages slowly: it is the slow adjustment
the optimum-currency-area model predicts.

\begin{record}[The euro, 1999 to today]
\begin{dated}{2026}
On 1 January 2026 Bulgaria introduced the euro, bringing the number of member
states of the European Union that use it to 21, with about 358 million people; the
lev was converted at 1.95583 per euro. % ledger: b3rg:ecb.bulgaria
\end{dated}
Every member gave up its exchange rate and its own interest rate, the two
instruments whose loss the optimum-currency-area criteria price.
\end{record}

\section{Currency boards and dollarisation}\label{sec:b3:exchange-rate-regimes:boards}

A \omterm{def:b3:exchange-rate-regimes:regimes}{currency board} removes the central bank's discretion: base money is issued only
against reserves, so it cannot finance deficits, and the first-generation attack
cannot happen. The cost is that the base shrinks whenever capital leaves, with no
lender of last resort to banks, and that the \omterm{def:b3:exchange-rate-regimes:regimes}{peg} can be undone only by changing
the law. \omterm{def:b3:exchange-rate-regimes:regimes}{Dollarisation} removes even that exit.

\begin{record}[Argentina's convertibility, 1991--2002]
The convertibility regime began with the law of 27 March 1991. % ledger: b3rg:arg.law
The official rate was 1.00 peso per dollar every year from 1995 to 2001, after
inflation that had seen the rate multiply by more than ten between 1989 and 1990.
Output fell in 1999, 2000 and 2001 ($-3.4$, $-0.8$ and $-4.4~\%$) with the \omterm{def:b3:exchange-rate-regimes:regimes}{peg} intact; the
law's core articles were repealed by a law published on 7 January 2002, the rate
averaged 3.06 pesos per dollar that year, and output fell by 10.9~\%. (Rates and
growth: World Bank, World Development Indicators.)
\end{record}

\section{Exercises}\label{sec:b3:exchange-rate-regimes:exercises}

\begin{exercise}[$\star$]\label{exo:b3:exchange-rate-regimes:1}
Name the eras visible in the opening chart and their dates. When was the share of
pegged economies highest?
\end{exercise}

\begin{exercise}[$\star$]\label{exo:b3:exchange-rate-regimes:2}
Which two corners of the trilemma does each choose: (a) a euro-area member; (b) a
floating economy with open capital markets; (c) a Bretton Woods \omterm{def:b3:exchange-rate-regimes:regimes}{peg} with capital
controls; (d) a country on the \omterm{def:b3:exchange-rate-regimes:gold-standard}{gold standard} in 1900?
\end{exercise}

\begin{exercise}[$\star$]\label{exo:b3:exchange-rate-regimes:3}
In the first-generation model with $k = 50$, $c = 100$, $\mu = 2$ and domestic credit of 40,
find the shadow exchange rate. If the \omterm{def:b3:exchange-rate-regimes:regimes}{peg} is 1.0, is the \omterm{def:b3:exchange-rate-regimes:regimes}{peg} under attack?
\end{exercise}

\begin{exercise}[$\star$]\label{exo:b3:exchange-rate-regimes:4}
What observation about a credible \omterm{def:b3:exchange-rate-regimes:regimes}{peg} with free capital movements would
contradict the trilemma?
\end{exercise}

\begin{exercise}[$\star\star$]\label{exo:b3:exchange-rate-regimes:5}
Estimate the trilemma regression on the JST data for \omterm{def:b3:exchange-rate-regimes:regimes}{pegs} and floats in four eras:
1870--1913, 1950--1972, 1973--1998 and 1999--2020. Where does the result follow the
trilemma, and where does it need the third corner to make sense?
\end{exercise}

\begin{exercise}[$\star\star$]\label{exo:b3:exchange-rate-regimes:6}
With $k = 100$, $c = 200$, a \omterm{def:b3:exchange-rate-regimes:regimes}{peg} of 1.0, reserves of 40 and credit growing at 4 a year,
find the date of the attack and the reserves lost in it. How does the date
change if $c$ doubles, and why?
\end{exercise}

\begin{exercise}[$\star\star$]\label{exo:b3:exchange-rate-regimes:7}
In the game of the chapter, the central bank's reserves are 14, then 5. Find the
equilibria in each case. Which policies move a country out of the middle range?
\end{exercise}

\begin{exercise}[$\star\star$]\label{exo:b3:exchange-rate-regimes:8}
From Eurostat's harmonised index, compute inflation in Greece and Germany each year
of 1999--2008, the mean gap and the rise of Greek prices relative to German ones.
\end{exercise}

\begin{exercise}[$\star\star$]\label{exo:b3:exchange-rate-regimes:9}
What evidence would show that the criteria for an \omterm{def:b3:exchange-rate-regimes:oca}{optimum currency area} are
endogenous? What would show they are not?
\end{exercise}

\begin{exercise}[$\star\star\star$]\label{exo:b3:exchange-rate-regimes:10}
Prove that in the first-generation model the exchange rate does not jump at the
attack, and that an attack at any other date would leave speculators with a loss
or an unexploited gain.
\end{exercise}

\begin{exercise}[$\star\star\star$]\label{exo:b3:exchange-rate-regimes:11}
A member of a currency union needs a real depreciation of 15~\% against its
partners, whose wages grow by 2~\% a year. If its own wages can at most be frozen,
how long does the depreciation take? If they can fall by 1~\% a year? (Take prices
proportional to wages.)
\end{exercise}

\begin{exercise}[$\star\star\star$]\label{exo:b3:exchange-rate-regimes:12}
A \omterm{def:b3:exchange-rate-regimes:regimes}{currency board} has a monetary base of 100, fully backed by reserves, and banks
hold reserves of 20 against deposits of 200. Capital of 8 leaves the country, drawn
from bank deposits and converted at the board. Find the new base and, if banks
restore their reserve ratio, the fall in deposits. What would a central bank have
done?
\end{exercise}

\section{Problem: Brenland Pegs the Dalen to the Kest}\label{sec:b3:exchange-rate-regimes:problem}

\begin{problem}[{Weekend problem --- Brenland pegs the dalen to the kest}]\label{pb:b3:exchange-rate-regimes:peg}
Brenland \omterm{def:b3:exchange-rate-regimes:regimes}{pegs} the dalen at $\bar S = 1.20$ dalens per kest. Money demand is $M = kS - c\dot S$
with $k = 100$ and $c = 200$ (billions of dalens; $\dot S$ in dalens per kest per year). Domestic
credit is 90 billion dalens at year 0 and grows by 3 billion a year to finance the
budget. % figdata: attack_stats, attack

\medskip\noindent\textbf{Part I --- the peg and its arithmetic.}
\begin{enumerate}
\item Find the money supply under the \omterm{def:b3:exchange-rate-regimes:regimes}{peg} and the reserves at year 0.
\item Why do reserves fall when credit grows under a \omterm{def:b3:exchange-rate-regimes:regimes}{peg}? Write the reserves at year $t$.
\item When would reserves run out if no one attacked?
\item Find the shadow exchange rate at year 0. Is the \omterm{def:b3:exchange-rate-regimes:regimes}{peg} under pressure?
\item Write the shadow rate at year $t$.
\item What is the kest value of the reserves at year 0?
\end{enumerate}

\medskip\noindent\textbf{Part II --- the attack.}
\begin{enumerate}[resume]
\item Find the date of the attack.
\item How many reserves are lost in the attack?
\item Find the money supply just after the attack.
\item At what rate does the dalen depreciate after the attack, in percent a year at
first?
\item By how much does Brenland's interest rate jump at the attack?
\item Why does the exchange rate not jump at the attack?
\item If the budget deficit doubled, so that credit grew by 6 billion a year, when
would the attack come?
\end{enumerate}

\medskip\noindent\textbf{Part III --- a scare.} At year 4, two funds can each sell 10 billion
dalens for kests, at a cost of 0.4 billion each; if the \omterm{def:b3:exchange-rate-regimes:regimes}{peg} falls, the dalen is
devalued by 20~\%.
\begin{enumerate}[resume]
\item Find the reserves at year 4, and the shadow rate. Do the fundamentals call for
an attack?
\item Write the payoff table of the two funds.
\item Find the Nash equilibria.
\item Unemployment is high, and the government will not raise rates to defend the
\omterm{def:b3:exchange-rate-regimes:regimes}{peg} once reserves fall below 12 billion. Find the equilibria now.
\item In what sense is an attack in Question 16 self-fulfilling? What does the
government's threshold in Question 17 change?
\item Name two things the government could have done to rule out the attack.
\end{enumerate}

\medskip\noindent\textbf{Part IV --- a currency board instead.}
\begin{enumerate}[resume]
\item What reserves would a \omterm{def:b3:exchange-rate-regimes:regimes}{currency board} need at year 0?
\item What share of them does Brenland hold? How large is the shortfall, in kests?
\item What happens to the budget's financing under the board?
\item Capital of 12 billion dalens leaves. Find the new base, and say what happens to
interest rates.
\item State the result: the date of the speculative attack, and the reserve cover a
\omterm{def:b3:exchange-rate-regimes:regimes}{currency board} needs.
\end{enumerate}
\end{problem}
